(1) 取入射单位Jones矢量为 $\boldsymbol e_{\rm in}=(1,1)^T/\sqrt2$；理想无损延迟片两轴的相对相位为 $\delta$，去除整体公共相位后出射为 $\boldsymbol e_{\rm out}=(1,e^{i\delta})^T/\sqrt2$。两分量强度各占一半而相位可不同，所以一般是椭圆偏振；只用光强标量无法保留所需相位信息。
(2) 当 $\delta=0$ 时两轴同相，光仍为沿 $45^\circ$ 的线偏振；$\delta=\pi/2$ 时两正交电场分量相差四分之一周期且振幅相等，电场端点沿圆转动，为圆偏振；$\delta=\pi$ 时两分量反相，出射线偏振沿 $-45^\circ$，与入射方向正交。圆偏振的左右旋标记取决于时间依赖与观察方向的约定。
(3) 与入射方向平行的分析器单位矢量为 $\boldsymbol a_+=(1,1)^T/\sqrt2$，透过复振幅为 $\boldsymbol a_+^\dagger\boldsymbol e_{\rm out}=(1+e^{i\delta})/2$。取模平方得 $I_+/I_0=|1+e^{i\delta}|^2/4=\cos^2(\delta/2)$。所以零延迟全透、四分之一波延迟透过一半、半波延迟完全消光，与前问偏振态一致。
(4) 垂直分析器矢量为 $\boldsymbol a_-=(1,-1)^T/\sqrt2$，相同计算得 $I_-/I_0=|1-e^{i\delta}|^2/4=\sin^2(\delta/2)$，并有 $I_++I_-=I_0$。对于左右旋圆偏振 $\delta=\pm\pi/2$，上述两个线检偏强度都为一半，无法判断旋向；必须再引入已知四分之一波片或测量包含相位信息的Stokes参量。
